\chapter{Running Example}\label{chap:running_example}

The simplest form of specialization using the tool is the following.

The default implementation does not need any change compared to standard Rust code.
The specialized implementation instead, is annotated with the attribute \inlinerust{#[when(...)]}, that lets us specify the conditions under which that implementation should be used.
In this case, we want to specialize the implementation of \texttt{MyTrait} for the type \texttt{i32}, so we write \inlinerust{#[when(T = i32)]}.
This will create a specialized version of the \texttt{trait} and of the \texttt{impl} for \texttt{i32} through a process that we are going to call \textit{monomorphic specialization}.

After doing this, we also need to instruct the method calls to use the correct version of the method.
This is done by using the \inlinerust{spec!} macro, which takes as first argument the method call and as second argument the type to specialize for.
As better explained in Chapter~\ref{chap:approach}, this is required because in the macro expansion phase the compiler doesn't give us enough information to automatically determine the types of the variables.

The complete code is shown in Listing~\ref{fig:monomorphic_specialization_pre}, where you can see the macros in action.
As you can probably understand, this approach preserves the soundness of the type system, since after the monomorphization process (see Listing~\ref{fig:monomorphic_specialization_post}) we have two completely separate versions of the trait and its implementation, that the compiler can easily distinguish and type check independently.

\begin{figure}[htb!]
	\centering
	\showrust*[0.98\linewidth]{monomorphized-specialization-pre.rs}
	\captionsetup{type=listing,skip=2pt}
	\caption{Example of the tool usage}
	\label{fig:monomorphic_specialization_pre}
\end{figure}

\begin{figure}[htb!]
	\centering
	\showrust*[0.98\linewidth]{monomorphized-specialization-post.rs}
	\captionsetup{type=listing,skip=2pt}
	\caption{Example of the code after our macros' expansion}
	\label{fig:monomorphic_specialization_post}
\end{figure}